\documentclass[aspectratio=169,10pt,english]{beamer}

\input{../../../_preambulo_beamer.tex}
\input{../../../_comandos_beamer.tex}

\selectlanguage{english}

\title{MA1001B - Statistical Modeling for Decision Making}
\subtitle{Lesson 19: Continuous Random Variables: Density and Interval Probability}
\author{}
\institute{Tecnol\'ogico de Monterrey}
\date{}

\begin{document}

\begin{frame}[plain]
  \vspace{1.2cm}
  {\Large\bfseries MA1001B - Statistical Modeling for Decision Making\par}
  \vspace{0.7cm}
  {\large Lesson 19: Continuous Random Variables\par}
  \vspace{0.7cm}
  {\color{TecRojo}\rule{\textwidth}{1.2pt}}
  \vspace{0.8cm}

  MA1001B Faculty\\[0.3cm]
  Term\\[0.3cm]
  Tecnol\'ogico de Monterrey
\end{frame}

\begin{frame}{The Continuous Probability Question}
  \begin{alertblock}{Guiding question}
    If service time can take any value on $[2, 14]$ minutes, where does
    probability live: at a point, or on an interval?
  \end{alertblock}

  \vspace{0.6em}
  By the end of this lesson, you can:
  \begin{itemize}
    \item distinguish a density $f(x)$ from a discrete PMF;
    \item read $P(a<X<b)$ as area under $f$;
    \item compute that area from $F(b)-F(a)$;
    \item explain why $P(X=c)=0$ for continuous $X$.
  \end{itemize}

  \vfill
  \scriptsize All service times used today are synthetic. No real company data are used.
\end{frame}

\begin{frame}{A Triangular Service-Time Density}
  \small
  Support $X\in[2,14]$ minutes, mode $6$. The peak height is
  \[
    f(6)=\frac{2}{14-2}=\frac{1}{6}=0.166667.
  \]
  The mean of a triangular density is
  \[
    E[X]=\frac{2+14+6}{3}=7.333333.
  \]

  \vspace{0.4em}
  \begin{exampleblock}{Density versus probability}
    $f(6)$ is a height. Probability is the area under $f$. The total area
    on $[2,14]$ is $1$.
  \end{exampleblock}
\end{frame}

\begin{frame}[fragile]{Step 1: Density Is a Height}
  \begin{columns}[T]
    \begin{column}{0.42\textwidth}
      \small
      \begin{lstlisting}
model = triang(
  c=1/3, loc=2, scale=12
)
f_mode = model.pdf(6)
area = model.cdf(14) - model.cdf(2)
      \end{lstlisting}

      \begin{block}{Verified output}
        $f(2)=0.000000$\\
        $f(6)=0.166667$\\
        Mean $=7.333333$\\
        Total area $=1.000000$
      \end{block}
    \end{column}
    \begin{column}{0.58\textwidth}
      \centering
      \includegraphics[width=\textwidth]{../figures/continuous_rv_01_density.png}
      \vspace{0.2em}
      \scriptsize Triangular density from Step 1
    \end{column}
  \end{columns}
\end{frame}

\begin{frame}{Interval Probability Is Area}
  \small
  For a continuous random variable,
  \[
    P(a<X<b)=F(b)-F(a)=\int_a^b f(x)\,dx.
  \]
  On the synthetic clock,
  \[
    F(4)=0.083333,\qquad F(10)=0.833333,
  \]
  \[
    P(4<X<10)=0.833333-0.083333=0.750000.
  \]
  The closed interval $P(4\le X\le 10)$ returns the same value.
\end{frame}

\begin{frame}[fragile]{Step 2: Shade $P(4<X<10)$}
  \begin{columns}[T]
    \begin{column}{0.40\textwidth}
      \small
      \begin{lstlisting}
area = (
    model.cdf(10)
    - model.cdf(4)
)
      \end{lstlisting}

      \begin{block}{Verified output}
        $F(4)=0.083333$\\
        $F(10)=0.833333$\\
        $P(4<X<10)=0.750000$\\
        $P(4\le X\le 10)=0.750000$
      \end{block}
    \end{column}
    \begin{column}{0.60\textwidth}
      \centering
      \includegraphics[width=\textwidth]{../figures/continuous_rv_02_interval_probability.png}
      \vspace{0.2em}
      \scriptsize Shaded interval from Step 2
    \end{column}
  \end{columns}
\end{frame}

\begin{frame}{Step 3: A Point Carries Probability 0}
  \begin{columns}[T]
    \begin{column}{0.42\textwidth}
      \small
      $P(|X-6|<0.50)=0.158854$\\
      $P(|X-6|<0.10)=0.033021$\\
      $P(|X-6|<0.01)=0.003330$

      \vspace{0.4em}
      \[
        F(6)=0.333333,\quad P(X=6)=0.
      \]

      \vspace{0.4em}
      \textbf{Limit:} the CDF has no jump, so a single minute has mass 0.
      Lesson 20 uses a flat uniform density and Monte Carlo.
    \end{column}
    \begin{column}{0.58\textwidth}
      \centering
      \includegraphics[width=\textwidth]{../figures/continuous_rv_03_point_mass_limit.png}
      \vspace{0.2em}
      \scriptsize CDF and shrinking windows from Step 3
    \end{column}
  \end{columns}
\end{frame}

\begin{frame}{Concept Check}
  \begin{enumerate}
    \item Why is $f(6)=0.166667$ not equal to $P(X=6)$?
    \item Compute $P(4<X<10)$ from $F(10)=0.833333$ and $F(4)=0.083333$.
    \item Why do $P(4<X<10)$ and $P(4\le X\le 10)$ match?
    \item What happens to $P(|X-6|<h)$ as $h\to 0$?
  \end{enumerate}

  \vspace{0.6em}
  \begin{block}{Connection to Lesson 20}
    The session can continue directly into the continuous uniform model and
    a Monte Carlo check of interval probability.
  \end{block}

  \vfill
  \scriptsize Source: OpenStax, \textit{Introductory Business Statistics 2e},
  Chapter 5, Section 5.1. CC BY-NC-SA.
\end{frame}

\appendix



\end{document}
